% 第 10 章　动量：碰撞的记账法
\chapter{动量：碰撞的记账法}

\step{反常识开场}

\begin{mainline}
厨房里有两只一模一样的鸡蛋，从同样的高度松手落下。一只落在瓷砖台面上，啪，碎了；另一只落在洗碗用的厚海绵上，弹了一下，完好无损。

这好像没什么可奇怪的——海绵软嘛。可是请把问题问得再细一点：两只鸡蛋从同一高度落下，落地前一瞬间的速度是一样的；落地之后，它们最终都停了下来。也就是说，两只鸡蛋经历的“从运动到静止”的过程，起点一样，终点也一样。瓷砖并没有让鸡蛋“停得更狠”——两只鸡蛋最后都是静止的，速度的改变量分毫不差。

那么，到底是什么不一样？是什么东西在瓷砖那一次“下手更重”？

再看一个反过来的例子。空手道表演者一掌劈断一块砖。如果他不用劈的，而是把手掌慢慢压上去，用尽全身力气，砖纹丝不动。同一只手，同一块砖，快和慢的区别为什么这么大？

还有一件事，可能你在台球桌或台球视频里见过：白球正正地撞上一颗静止的球，撞完之后，白球停在原地，那颗球却用几乎相同的速度滚走了——好像白球把自己的“运动”整个过户给了对方。运动这种东西，难道可以像快递包裹一样，从一个物体手里转交到另一个物体手里，还当面点清、一件不少？

这一章要回答的就是这些问题。答案藏在一个新的物理量里。它一旦登场，你会发现：碰撞、爆炸、火箭、后坐力，这些看起来毫不相干的事情，其实都在记同一本账。
\end{mainline}

\step{先猜一猜}

照例，先押三注。请把答案写在纸上，本章结尾对账。

第一题。两只相同的鸡蛋从同一高度落下，一只落在瓷砖上，一只落在海绵上，都停住了。哪只鸡蛋的“速度改变量”更大？（a）瓷砖上的那只，因为它停得更猛；（b）海绵上的那只，因为它还弹了一下；（c）一样大。

第二题。一辆空卡车和一辆小轿车迎面相撞。撞的一瞬间，谁受到的冲击力更大？（a）小轿车，它轻，吃亏；（b）卡车，它结实，顶得住；（c）一样大。

第三题。一颗静止悬空的鞭炮炸成两半，一半向左飞，一半向右飞。哪一半飞得更快？（a）重的那半；（b）轻的那半；（c）一样快；（d）不一定，看火药怎么分布。

先别往下看，给出你的答案，顺便写一句理由。理由比答案重要。

\step{动手实验}

\begin{experiment}[桌面上的碰撞快递]
\textbf{材料}：七八枚相同的硬币（一元硬币最好）、一把直尺、一张光滑的桌面、一颗生鸡蛋、一块厚海绵（或叠几层的毛巾）、一部手机（用来慢动作录像，可选）。

\textbf{第一部分：撞出几枚？}把四枚硬币在桌面上排成一条直线，互相贴紧。把直尺放在这排硬币的正后方，让第五枚硬币贴着直尺边缘滑过去——直尺保证它走直线——正正地撞在这排硬币的队尾上。观察结果：最前面的\emph{一枚}硬币弹了出去，其余几乎不动，撞进来的那枚停了下来。现在用两枚硬币并排一起撞进去：这回弹出来的是\emph{两枚}。试一试三枚。你会发现一条朴素的规律：进去几枚，出来几枚，而且出去的硬币大约带走了进来的硬币的全部“运动”。排中间的那些硬币仿佛只是“转手”了一下。

\textbf{第二部分：鸡蛋的两次落地。}把海绵铺在台面上，把生鸡蛋从大约三十厘米高处松手，落在海绵上：完好（多做两次，放心，海绵很可靠）。然后把鸡蛋从\emph{同样}的高度落在坚硬的台面上——如果你舍不得浪费一颗鸡蛋，可以把这一步换成慢动作录像：用手机拍一只装了水的薄塑料袋，从同一高度分别落到海绵和台面上，回看录像。注意看两件事：袋子（或鸡蛋）从开始接触底面、到完全停下来，各花了多少时间。你会发现，海绵上的那一次，停下来的过程明显拖得更长。

\textbf{想一想}：第一部分里，“进去几枚出来几枚”像是某种东西被严格地转手了，那是什么东西？为什么不是“速度”本身——毕竟两枚一起撞进去时，弹出来的两枚每一枚都比单枚撞击时飞得慢？第二部分里，两次落下的速度一样、最后都停住，不同的只有停下来的\emph{时间}。时间和结果之间，藏着什么关系？
\end{experiment}

\step{旧模型遇到麻烦}

到目前为止，我们手里的运动学工具是第 5、6 章给的：位置、速度、加速度，外加牛顿第二定律——合力决定加速度。这套工具在斜坡、抛物、圆周运动里所向披靡，可是一碰到“碰撞”，它立刻显出三个软肋。

第一，\emph{碰撞的力我们根本不知道}。鸡蛋撞上瓷砖的那零点几秒里，瓷砖给鸡蛋的力是多大？它随时间怎么变？一开始接触时力多大，蛋壳开始出现裂纹时力多大？没有人知道，也没有必要知道——这个力巨大、短暂、形状复杂，连测量都很困难。牛顿第二定律要你提供每个时刻的力，才能算出每个时刻的运动；可碰撞偏偏把这个关键输入藏了起来。

第二，\emph{能量守恒不够数}。第 9 章给了我们能量这本账。试试看能不能用能量守恒解决台球问题：白球以速度 \(v\) 撞上静止的、同样质量的球，撞完两球各是什么速度？能量守恒只给你一个方程，可未知数有两个——两颗球各自的速度。一个方程解不出两个未知数。更糟的是，大量碰撞里动能根本不守恒：两团橡皮泥撞在一起粘住了，动能凭空少了一大截（变成了热和形变），能量总账当然要平，可动能账已经亏了。光有能量守恒，我们连“撞完大概什么样”都说不准。

第三，\emph{有些过程连“力”都懒得出场}。鞭炮爆炸，火药的内力把两半推开；人在小船头上往船尾走，船往相反方向挪；开枪时枪托猛地向后撞你的肩膀。这些过程里的力全是“内部事务”，细节一团乱麻，可结果却整齐得惊人——总有什么东西前后不变。

三个软肋指向同一个缺口：我们需要一个新的物理量，它不关心碰撞过程中每一毫秒的细节，只认\emph{碰撞前后}的总数；它在物体之间转手时，总量一分不多、一分不少。一句话，我们需要一本新的账。

\step{发明新概念}

现在，像历史上的物理学家那样，把这本账发明出来。

先说实验里那条线索。桌面上“进去几枚出来几枚”，说明被转手的东西跟硬币的\emph{枚数}成正比——也就是跟质量成正比。而你又注意到，两枚撞进去时，弹出的每枚比单枚撞击时慢——说明被转手的总量是固定的，摊到两枚上，每枚分到的速度就薄了。把这两条拼起来：被转手的东西，既含质量，又含速度，而且大致是\emph{质量乘速度}。给它起个名字，叫\emph{动量}：
\[
p=mv .
\]
动量有方向，和速度同向。一个物体动量大，可能是它重，也可能是它快——慢慢开的大卡车和飞驰的子弹，动量可以不相上下。这正是“又重又快的东西难拦住”的精确说法。

再看鸡蛋的线索。两次落地，速度从同一个 \(v\) 变成 \(0\)，动量的改变量 \(mv\) 一模一样；不同的是海绵把停下来的时间拖长了。瓷砖那次，同样的动量改变挤在极短的时间里完成；海绵那次，摊在较长的时间里完成。于是猜：\emph{力，就是单位时间里改变的动量}。同样的动量改变，时间越短，力越狠。写成式子：
\[
F=\frac{\Delta p}{\Delta t},\qquad\text{或者}\qquad F\,\Delta t=\Delta p .
\]
这不是一条新定律，它就是把第 6 章的 \(F=ma\) 换个写法（质量不变时，\(\Delta p=m\Delta v\)，两边除以 \(\Delta t\) 就是 \(F=ma\)）。但换个写法，世界忽然亮了一块：左边 \(F\Delta t\) 叫\emph{冲量}——力在时间上的积累；右边是动量的改变。这个形式告诉我们一件 \(F=ma\) 藏得很深的事：\emph{力的作用效果，要看它坚持了多久}。轻拍一下和猛推一把，即使力的峰值一样，坚持的时间不同，交出去的动量就不同。

有了动量，碰撞的秘密就剩一层窗户纸了。看两个物体 \(A\) 和 \(B\) 相撞。第 6 章的第三条规定：\(B\) 给 \(A\) 的力和 \(A\) 给 \(B\) 的力，大小相等、方向相反，而且\emph{同时存在、同时消失}——持续的时间当然也是同一段 \(\Delta t\)。于是这段时间里，\(A\) 收到的冲量和 \(B\) 收到的冲量等大反向，\(A\) 增加的动量恰好等于 \(B\) 减少的动量。总量呢？一分没动：
\[
p_{A}+p_{B}=\text{常数}.
\]
这就是\emph{动量守恒}。注意我们几乎什么都没假设：没假设力是弹性的还是塑性的，没假设撞完粘不粘，甚至没假设我们知道力是多大——只用了“相互作用成对出现、同时同长”。这就是为什么动量守恒对碰撞的细节一问三不知却百发百中：它根本不在乎细节。

但守恒是有条件的，条件就写在推导里：我们只管了 \(A\) 和 \(B\) 之间的力。如果还有\emph{第三方}插手——比如地面摩擦、墙壁顶着、重力在旁边猛拽——那这些外力的冲量也会往账里添进项或减项，总动量就不守恒了。所以准确的说法是：\emph{一个系统，当外力对它的总冲量为零（或在关心的过程中可以忽略）时，它的总动量守恒}。这里第一次要认真对待“系统”这个词：把边界划在哪，决定了哪些力算内力（只在账内转手，不改总量），哪些力算外力（会直接改动总账）。划系统不是文字游戏，是动量记账的第一步。

最后补一个看整体的新视角。既然总动量守恒，那“总动量除以总质量”这个速度也守恒。它是什么东西的速度？对两个质量分别为 \(m_1,m_2\)、位置在 \(x_1,x_2\) 的物体，定义
\[
x_{\text{质心}}=\frac{m_1x_1+m_2x_2}{m_1+m_2},
\]
叫系统的\emph{质心}——按质量加权的平均位置，重物把它往自己那边拽。可以验证（数学桥层里做）：质心的速度正好是“总动量除以总质量”。于是动量守恒有一句等价的话：\emph{不管内部怎么撞、怎么炸，只要外力冲量为零，质心就自顾自地匀速直线运动}。烟花在最高点炸开，碎片四散飞射，乱成一锅粥——可所有碎片的质心，依然沿着没炸时那条抛物线继续走，好像爆炸从未发生。

\step{一句话模型}

\begin{oneline}
动量 \(p=mv\) 是“运动量”的账：力是动量的搬运工，冲量 \(F\Delta t\) 是搬运的凭据；系统内物体之间撞来撞去、炸来炸去，只是把动量在内部转手——只要外力不插手，总账一分不多一分不少。
\end{oneline}

\step{数学翻译器}

把账写成式子，每条写完立刻用极端情形检查。

\textbf{动量：\(p=mv\)。}描述什么：一个物体带着多少“运动量”，朝哪个方向。为什么这样写：桌面硬币实验说它被转手时既正比于质量、又正比于速度，乘法是最简单的满足两条的写法。检查：\(v=0\) 时 \(p=0\)，静止的东西没有运动量，对。\(m\) 翻倍 \(p\) 翻倍，两枚硬币撞进去就要吐出两枚，对。方向反转：速度取负（掉头），动量取负，正负相消——这就是为什么向左的和向右的动量可以互相抵消，爆炸后两边合起来仍然为零。极大检查：一列万吨列车以步行速度开动，\(p=10^{7}\times1=10^{7}\) 千克米每秒，和一辆轿车以两千倍音速飞行相当——这就是为什么看起来“慢慢挪”的火车头照样能把汽车推平。

\textbf{冲量与动量定理：\(F\Delta t=\Delta p\)。}描述什么：一段力的作用，给物体的动量账户存进（或取出）多少。为什么这样写：它就是把 \(F=ma\) 两边乘以时间重排的结果，不添新物理，只换视角。何时成立：力恒定时直接成立；力变化时，要用每个瞬间的力累加（见数学桥层）。何时失效：它本身几乎不失效，但用它估算时若把时间猜错，结果会离谱。检查：\(\Delta t\) 极小而 \(\Delta p\) 固定，则 \(F\) 冲向巨大——这就是瓷砖杀鸡蛋的数学画像；同一个 \(\Delta p\) 摊到很长的 \(\Delta t\) 上，\(F\) 就温柔——这是海绵的画像。冲量还可以画成图：横轴时间、纵轴力，曲线下的面积就是冲量。两只鸡蛋的两条曲线形状天差地别，面积却相等。

\begin{center}
\begin{tikzpicture}[>=Stealth,scale=0.9]
  % 左图：硬着陆
  \draw[->] (0,0) -- (4.6,0) node[right]{$t$};
  \draw[->] (0,0) -- (0,3.0) node[above]{$F$};
  \draw[very thick,red!70!black] (0.4,0) .. controls (0.7,2.6) and (0.9,2.6) .. (1.2,0);
  \fill[red!15] (0.4,0) .. controls (0.7,2.6) and (0.9,2.6) .. (1.2,0) -- cycle;
  \node at (0.8,-0.5) {瓷砖：窄而高};
  % 右图：软着陆
  \begin{scope}[xshift=6.2cm]
  \draw[->] (0,0) -- (4.6,0) node[right]{$t$};
  \draw[->] (0,0) -- (0,3.0) node[above]{$F$};
  \draw[very thick,blue!70!black] (0.3,0) .. controls (1.3,1.0) and (2.7,1.0) .. (3.7,0);
  \fill[blue!15] (0.3,0) .. controls (1.3,1.0) and (2.7,1.0) .. (3.7,0) -- cycle;
  \node at (2.0,-0.5) {海绵：宽而矮};
  \end{scope}
\end{tikzpicture}
\end{center}

\textbf{一维碰撞的守恒式。}两个物体沿一条线运动，碰前速度 \(v_1,v_2\)，碰后 \(v_1',v_2'\)，动量守恒写作
\[
m_1v_1+m_2v_2=m_1v_1'+m_2v_2' .
\]
注意这是一个\emph{带正负号}的式子：朝左的速度记负，式子自动处理方向。

\emph{弹性碰撞}（动能也守恒，比如台球、硬币那种硬碰硬）有个漂亮推论，推导放在数学桥层：
\[
v_1'=\frac{m_1-m_2}{m_1+m_2}v_1+\frac{2m_2}{m_1+m_2}v_2,\qquad
v_2'=\frac{2m_1}{m_1+m_2}v_1+\frac{m_2-m_1}{m_1+m_2}v_2 .
\]
检查三个特殊情形。其一，等质量 \(m_1=m_2\)：式子变成 \(v_1'=v_2,\ v_2'=v_1\)——\emph{交换速度}。这就是白球撞停自己、目标球接走全部速度的原因，也是“进去几枚出来几枚”的算术根源。其二，轻的撞不动的墙（\(m_1\ll m_2\)，\(v_2=0\)）：\(v_1'\approx -v_1\)，原速弹回；墙几乎不动，但拿走了 \(2m_1v_1\) 的动量——皮球砸墙，球的速度只掉了个头，动量却变了 \(2mv\)，差额全给了墙（和墙连着的地球）。其三，也是最反直觉的：\emph{重的撞轻的}（\(m_1\gg m_2\)，\(v_2=0\)）：\(v_2'\approx 2v_1\)——静止的轻物被撞后飞出的速度，可以是撞入者速度的\emph{两倍}！直觉总说“被撞的顶多和撞的一样快”，式子说不。保龄球撞飞球瓶、大锤敲击让小钉子飞溅，都有它的影子。

\emph{完全非弹性碰撞}（撞完粘在一起，比如两团橡皮泥、子弹钻进木块）：碰后共速，动量守恒直接给出
\[
v'=\frac{m_1v_1+m_2v_2}{m_1+m_2}.
\]
检查：\(m_2\) 原本静止且质量巨大，则 \(v'\) 很小——子弹钻进大地，大地几乎不动，但动量一分不少地收下了。现在算动能账：碰前动能 \(\tfrac12 m_1v_1^2\)，碰后 \(\tfrac12(m_1+m_2)v'^2\)，代入可知碰后\emph{必然更小}，亏空变成了热、声和永久形变。这里是本章最重要的一个对照：\emph{动量守恒和动能守恒是两本独立的账}。碰撞中动量总是守恒（外力可忽略时），动能却可以大量蒸发；反过来，炸药爆炸时动能凭空增加（化学能变成动能），动量却依然守恒——炸前为零，炸后各碎片的动量加起来还是零。想用能量守恒代替动量守恒去解碰撞，就像只拿着收入账去对支出账，注定对不平。

\textbf{带真实数据的估算。}一辆轿车以 60 千米每小时（约 16.7 米每秒）撞墙，驾驶员质量 60 千克，系着安全带随车一起停下。驾驶员的动量改变量 \(\Delta p=60\times16.7\approx1000\) 千克米每秒。如果胸口直接撞上方向盘，在约 0.05 秒内停下：平均力 \(F\approx1000/0.05=20000\) 牛——相当于两吨重物的重量压在胸口，肋骨必断。如果安全气囊把停下过程拉长到约 0.3 秒：平均力降到约 3300 牛，疼，但活得下来。气囊没有减少一丁点动量改变，它只是\emph{把同一份冲量摊薄在更长的时间里}。安全带、头盔内衬、跳远的沙坑、接棒球时顺势后缩的手、猫落地时弯曲的腿，全都是同一条原理的应用。

\begin{mathbridge}
\textbf{变力的冲量。}力随时间变化时，把碰撞过程切成许多小段，每段里力近似恒定，冲量是 \(F_i\Delta t_i\)，全部加起来取极限，就是积分：
\[
J=\int_{t_1}^{t_2} F(t)\,dt=\Delta p .
\]
图上那条曲线下方的面积，严格意义就在这里。

\textbf{弹性碰撞公式的推导。}联立动量守恒 \(m_1v_1+m_2v_2=m_1v_1'+m_2v_2'\) 与动能守恒 \(\tfrac12 m_1v_1^2+\tfrac12 m_2v_2^2=\tfrac12 m_1v_1'^2+\tfrac12 m_2v_2'^2\)。把两式分别移项相除，可化简出一个极简的中间结论：\(v_2'-v_1'=v_1-v_2\)，即“分离速度等于接近速度”（它和动能守恒等价，但好算得多）。把它和动量守恒式联立，两个线性方程解两个未知数，就得到正文的结果。

\textbf{质心速度。}对 \(x_{\text{质心}}=(m_1x_1+m_2x_2)/(m_1+m_2)\) 两边取时间变化率：
\[
v_{\text{质心}}=\frac{m_1v_1+m_2v_2}{m_1+m_2}=\frac{p_{\text{总}}}{m_{\text{总}}}.
\]
总动量守恒，等价于质心匀速直线运动。外力的总冲量改变的，也正是“质心的运动”这一样东西——不管系统内部是两个人、一万块碎片还是一团气体。多个物体时把求和项数加长即可，写成 \(x_{\text{质心}}=\sum_i m_ix_i/\sum_i m_i\)。
\end{mathbridge}

\step{模型的边界}

动量守恒强大，但有三道边界，跨过去就会出错。

第一，\emph{外力的账不能不认}。动量守恒的前提是系统所受外力的总冲量为零或可以忽略。台球桌上，击球之前桌面摩擦一直在悄悄抽税；球撞库边时，库边连着球桌、球桌连着地球，动量流进了大地，只是地球质量太大，看不出动静。什么时候可以“假装”外力不存在？碰撞时间极短时：撞的那零点几秒里，摩擦这类常规力的冲量 \(F\Delta t\) 小得可以忽略，于是“碰撞前后那一瞬间”动量近似守恒——这就是为什么子弹打木块这类问题可以先守恒、再考虑摩擦。

第二，\emph{动量是带方向的}。总动量不守恒，不代表每个方向都不守恒。斜着抛出又炸开的烟花，竖直方向有重力在不断注入冲量，竖直动量不守恒；但水平方向没有外力，水平动量照样守恒，质心照样走它水平匀速的那一份。分方向记账，是把“矢量”两个字落到实处的第一功。

第三，\emph{经典动量有速度上限}。当速度接近光速，\(p=mv\) 这本账记不下了——以恒定大小的力推一个物体，按 \(F\Delta t=m\Delta v\) 它的速度应该无限增长，可实验发现加速越来越难，速度始终够不着光速。不是动量守恒错了，而是动量的\emph{定义}需要升级：高速世界里，动量和速度之间不再是简单的正比关系（本书采用的说法是：物体的质量本身不变，变的是动量与速度的关系，详见第 40 章）。在粒子加速器里，物理学家每天用升级后的动量账本核对粒子碰撞，账依然分毫不差——动量守恒本身，比牛顿力学的适用范围活得久得多。

最后点名两个典型误用。其一，把动量当成“物体储存的力气”：说子弹“带着很大的力”打穿木板——错，子弹带的是动量；力只在它和木板相互作用的那段时间里存在，大小取决于动量改变得有多快。其二，以为“动量守恒所以撞不坏”：两辆车对撞，总动量确实守恒，可每辆车\emph{自己}的动量改变照样巨大，该断的肋骨一根不少。守恒的是总账，不是平安。

\step{回头解释开场现象}

现在可以对账了。

两只鸡蛋，同一高度落下，落地前动量同为 \(mv\)，落地后同为零——动量改变量一模一样，冲量一模一样，图上两块面积一样大。瓷砖把这份冲量挤进几毫秒，峰值力冲破蛋壳的承受上限；海绵把同样的冲量摊到几十倍长的时间，峰值力降到蛋壳能扛住的范围。杀死鸡蛋的不是“停得猛”，而是“停得快”。（第一题答案：一样大。）

空手道劈砖是同一个原理反过来用：手掌带着动量高速拍下，砖被下面的支点顶住，手掌的动量必须在极短时间内交出去，于是峰值力极大——砖断了。慢慢压，同样的手，动量变化细水长流，力峰值连砖的零头都够不着。

台球白球“过户”运动，是等质量弹性碰撞的交换速度：白球交出全部动量，目标球接走，桌面上的硬币亦然。第二题答案是一样大——相互作用力成对出现；小轿车吃亏不是因为它受力大，而是同样的力让轻的它速度改变更猛：\(F\Delta t=m\Delta v\) 里 \(m\) 小，\(\Delta v\) 就大，变形和伤亡由此而来。

鞭炮两半，炸前总动量为零，炸后也必须为零：\(m_1v_1+m_2v_2=0\)，轻的那半速度大。（第三题答案：轻的那半快。）枪的后坐、你从船头跳上岸时船向后退，全是这一条。

\step{脑洞实验与挑战题}

\begin{thought}[太空里，火箭推什么？]
在地面上，船靠桨推水，车靠轮推地，人靠脚蹬地——前进的本质都是“把别的东西向后推，自己收反方向的动量”。可是火箭飞出大气层，四周真空，无桨可划、无地可蹬，它推什么？

答案是：它推自己带的东西。火箭把燃料烧成的燃气以每秒几千米的速度向后喷出，每一团燃气带走一份向后的动量，火箭就收下一份等大的向前的动量。火箭不是在“推真空”，而是在不停地\emph{向后扔东西}——只不过扔的是高温气体，扔得又快又狠。土星五号火箭起飞时重约 3000 吨，第一级发动机每秒烧掉约 13 吨燃料，靠的就是这条最原始的账。

由此还得到一个不好听但重要的结论：火箭的燃料不仅要推动卫星，还要推动\emph{还没烧掉的燃料自己}。想飞得快，就得多带燃料；多带燃料，又得更费力地推这些燃料。这就是为什么去趟月球要用三千吨的火箭送一个几十吨的舱段——动量账在太空里收税极狠（具体的定量关系在挑战层）。
\end{thought}

\begin{challenge}
\textbf{火箭方程。}设火箭某时刻质量为 \(m\)、速度为 \(v\)，在极短时间里向后喷出质量为 \((-dm)\) 的燃气，燃气相对火箭的喷射速度恒为 \(u\)。对“火箭加燃气”系统用动量守恒（太空里外力忽略），展开并略去二阶小量，得
\[
m\,dv=-u\,dm .
\]
从点火到燃料烧完积分，得齐奥尔科夫斯基火箭方程：
\[
\Delta v=u\ln\frac{m_0}{m_{\text{末}}} .
\]
它量化了正文里那句“收税极狠”：想让 \(\Delta v\) 翻倍，质量比 \(m_0/m_{\text{末}}\) 要变成原来的\emph{平方}——对数增长意味着燃料需求爆炸式增长。多级火箭把烧完的空壳半路扔掉，就是在和这条对数抢钱。

\textbf{光也有动量。}后面第 25、40 章会看到，连没有静质量的光也携带动量，照在物体上会产生微小的压力。太阳帆飞船正是打算靠这份光压，不带一克燃料在行星间航行——动量守恒管到了光头上。再往后第 58 章会揭晓这本账的最终来历：动量守恒源于“空间处处平等”——物理规律不因你平移一步而改变。本章学会的记账法，在那里会升级成物理学最深的一条原理。
\end{challenge}

\begin{puzzles}
\item 帆船漫画里有个著名桥段：船员在甲板上架一台大功率电风扇，对着自己的帆猛吹，船能前进吗？如果他把风扇转一百八十度、朝船尾后方的空气吹呢？\emph{提示}：把“船加风扇加帆”划成一个系统，风扇吹出的风和帆挡住的风，是系统内部事务还是外部事务？两种情形的边界划法一样吗？
\item 一辆敞篷平板货车在平直轨道上无摩擦地匀速滑行，突然下起大雨，雨点竖直落进敞开的车厢并积在里面。车速会变大、变小还是不变？雨停后车厢底部漏了个洞，积水慢慢漏光，车速又会怎样？\emph{提示}：雨点落下前水平速度是多少？它进入车厢后，水平动量从哪来？漏下去的水离开车厢时，水平速度是多少？
\item 一颗皮球垂直砸向一堵固定在地上的墙，原速弹回。动能守恒了，动量却不守恒——球的动量掉了个头，变了 \(2mv\)。这违反动量守恒吗？那笔动量去哪了？为什么墙（和地球）收下了动量却几乎不表现动能？\emph{提示}：动量是 \(mv\)，动能是 \(\tfrac12 mv^2\)；把同样一份动量交给一个质量大 \(10^{20}\) 倍的对象，它的速度和动能各缩成多少？
\item 你和一位体重是你两倍的朋友，面对面站在绝对光滑的冰面上，互相猛推一把后各自滑开。谁滑得快？滑出同样长时间后，谁离出发点远？你们俩的质心在哪、动没动？\emph{提示}：互推的力等大反向、同时同长，两边收到的冲量是什么关系？把坐标原点设在哪儿，会影响质心动不动这个结论吗？
\end{puzzles}
